\nonstopmode{}
\documentclass[a4paper]{book}
\usepackage[times,inconsolata,hyper]{Rd}
\usepackage{makeidx}
\makeatletter\@ifl@t@r\fmtversion{2018/04/01}{}{\usepackage[utf8]{inputenc}}\makeatother
% \usepackage{graphicx} % @USE GRAPHICX@
\makeindex{}
\begin{document}
\chapter*{}
\begin{center}
{\textbf{\huge Package `PALM'}}
\par\bigskip{\large \today}
\end{center}
\ifthenelse{\boolean{Rd@use@hyper}}{\hypersetup{pdftitle = {PALM: Fast and reliable association discovery in large-scale microbiome studies and meta-analyses}}}{}
\begin{description}
\raggedright{}
\item[Type]\AsIs{Package}
\item[Title]\AsIs{Fast and reliable association discovery in large-scale microbiome studies and meta-analyses}
\item[Version]\AsIs{0.1.0}
\item[Author]\AsIs{Zhoujingpeng Wei}
\item[Maintainer]\AsIs{Zhoujingpeng Wei }\email{zwei74@wisc.edu}\AsIs{}
\item[Description]\AsIs{This R package implements a quasi-Poisson-based framework designed for robust, scalable, and reproducible identification of covariate-associated microbial features in large-scale microbiome association studies and meta-analyses.}
\item[Data]\AsIs{2025-11-06}
\item[Reference]\AsIs{Wei et al. 'Fast and reliable association discovery in large-scale microbiome studies and meta-analyses using PALM'. Submitted.}
\item[License]\AsIs{GPL (>= 3) | file LICENSE}
\item[LinkingTo]\AsIs{Rcpp, RcppArmadillo}
\item[Imports]\AsIs{MASS (>= 7.3-60),
metafor (>= 4.8.0),
abess (>= 0.4.10),
dplyr (>= 1.1.4),
brglm2 (>= 0.9.2),
Rcpp (>= 1.0.14),
RcppArmadillo (>= 14.2.2-1)}
\item[Encoding]\AsIs{UTF-8}
\item[LazyData]\AsIs{true}
\item[RoxygenNote]\AsIs{7.3.2}
\item[Suggests]\AsIs{knitr,
rmarkdown}
\item[Depends]\AsIs{R (>= 4.3.0)}
\item[VignetteBuilder]\AsIs{knitr}
\end{description}
\Rdcontents{Contents}
\HeaderA{CRC\_data}{Datasets from two metagenomics studies of colorectal cancer (CRC)}{CRC.Rul.data}
\keyword{datasets}{CRC\_data}
%
\begin{Description}
The "CRC\_data" includes "CRC\_abd" and "CRC\_meta". "CRC\_abd" is a sample-by-feature matrix of
relative abundance counts from the two studies including 267 species under order "Clostridiales".
The "CRC\_meta" is a data frame including the sample-level variables from the two studies.
\end{Description}
%
\begin{Usage}
\begin{verbatim}
data(CRC_data)
\end{verbatim}
\end{Usage}
%
\begin{Format}
An object of class \code{list} of length 2.
\end{Format}
%
\begin{Source}
<https://github.com/zellerlab/crc\_meta>
\end{Source}
%
\begin{References}
Wirbel, Jakob et al. Nat Med. 2019 Apr;25(4):679-689.
\end{References}
%
\begin{Examples}
\begin{ExampleCode}

library("PALM")
data("CRC_data")
CRC_abd <- CRC_data$CRC_abd
CRC_meta <- CRC_data$CRC_meta


\end{ExampleCode}
\end{Examples}
\HeaderA{palm}{Meta-analysis for feature-level association testing for each covariate of interest}{palm}
%
\begin{Usage}
\begin{verbatim}
palm(
  rel.abd,
  covariate.interest,
  covariate.adjust = NULL,
  cluster = NULL,
  depth = NULL,
  depth.filter = 0,
  prev.filter = 0.1,
  p.adjust.method = "fdr",
  meta.method = "EE",
  correct = "median"
)
\end{verbatim}
\end{Usage}
%
\begin{Arguments}
\begin{ldescription}
\item[\code{rel.abd}] For a single association study, provide a matrix of relative abundance counts,
with sample IDs as row names and microbial feature IDs as column names.
The matrix must not contain any missing values.
For a meta-analysis, provide a list of such matrices, where each element corresponds
to one study, and the name of each element represents the study ID.

\item[\code{covariate.interest}] For a single association study, provide a matrix where rows represent samples
and columns represent numeric covariates of interest. The order of samples must match the order in \code{rel.abd},
and the matrix must not contain any missing values.
For a meta-analysis, provide a list of such matrices, where each element corresponds to one study,
and the name of each element represents the study ID. The set of covariates may differ across studies.

\item[\code{covariate.adjust}] For a single association study, provide a data frame where rows represent samples
and columns represent covariates to be adjusted for in the model. Covariates must be either factors or numeric vectors.
The order of samples must match the order in \code{rel.abd}, and the data frame must not contain missing values.
For a meta-analysis, provide a list of such data frames, where each element corresponds to one study,
and the name of each element represents the study ID. The set of covariates may differ across studies.
The default is \code{NULL}, meaning no covariates are adjusted.

\item[\code{cluster}] For a single association study with correlated samples, provide a vector defining
the sample clusters. The order of samples in the vector must match the order in \code{rel.abd}.
For example, the values may represent subject IDs if each subject has multiple correlated samples
(e.g., in a longitudinal study).
For a meta-analysis, provide a list of such vectors, where each element pertains to one study with correlated samples,
and the name of each element represents the study ID. The order of samples in each vector must match the order in
\code{rel.abd} for the corresponding study.
The default is \code{NULL}, assuming all samples across all studies are independent.

\item[\code{depth}] For a single association study, provide a vector representing the sequencing depth for each sample.
The length of the vector must match the number of samples, and the order must align with \code{rel.abd}.
For a meta-analysis, provide a list of such vectors, where each element corresponds to one study,
and the name of each element represents the study ID.
The default is \code{NULL}, in which case the sequencing depth is calculated as the row sums of \code{rel.abd}.

\item[\code{depth.filter}] A cutoff value used to remove samples with sequencing depth less than or equal to the cutoff.
Default is \code{0}.

\item[\code{prev.filter}] A cutoff value to remove microbial features with prevalence (proportion of nonzero observations)
less than or equal to the cutoff. This cutoff is applied to each study, so a feature may be removed in a subset of studies.
The cutoff value must be in the range \code{0–1}. Default is \code{0.1}.

\item[\code{p.adjust.method}] Character string specifying the multiple testing correction method.
Options include \code{"holm"}, \code{"hochberg"}, \code{"hommel"}, \code{"bonferroni"}, \code{"BH"},
\code{"BY"}, \code{"fdr"}, and \code{"none"}. See \code{?stats::p.adjust} for more details.

\item[\code{meta.method}] Character string specifying whether an equal-effects or random-effects model should be fitted.
An equal-effects model is fitted when \code{method = "EE"}.
A random-effects model is fitted by setting \code{method} to one of the following:
\code{"DL"}, \code{"HE"}, \code{"HS"}, \code{"HSk"}, \code{"SJ"}, \code{"ML"}, \code{"REML"},
\code{"EB"}, \code{"PM"}, \code{"GENQ"}, \code{"PMM"}, or \code{"GENQM"}.
The default is \code{"EE"}. See \code{?metafor::rma} for more details.

\item[\code{correct}] Compositional effect correction method. The default is \code{"median"}.
If \code{"median"}, the compositional effect is corrected using the median of the estimates.
If \code{"tune"}, the compositional effect is corrected using a data-driven tuned value.
If \code{NULL}, no compositional correction is applied (i.e., the analysis uses relative-abundance summary
statistics directly and outputs RA-level summary statistics).
Recommendation: For sparse signals (≤ 20
For moderately dense signals (> 20
Both methods assume that the proportion of differential AA features is < 50

\end{ldescription}

A list containing one component for each covariate of interest.

For single-study analysis, the component includes:


For meta-analysis across multiple studies, the component includes:



Implements the PALM framework for both single-study association analysis and
meta-analysis across multiple studies. This function conducts feature-level
association testing, evaluating one covariate of interest at a time.



library(PALM)
data("CRC\_data", package = "PALM")
CRC\_abd <- CRC\_data\$CRC\_abd
CRC\_meta <- CRC\_data\$CRC\_meta

\#\#\#\#\#\#\#\#\#\# Generate summary statistics \#\#\#\#\#\#\#\#\#\#
rel.abd <- list()
covariate.interest <- list()
for (d in unique(CRC\_meta\$Study)) \{
  rel.abd[[d]] <- CRC\_abd[CRC\_meta\$Sample\_ID[CRC\_meta\$Study == d], ]
  disease <- as.numeric(CRC\_meta\$Group[CRC\_meta\$Study == d] == "CRC")
  names(disease) <- CRC\_meta\$Sample\_ID[CRC\_meta\$Study == d]
  covariate.interest[[d]] <- matrix(disease, ncol = 1,
                                    dimnames = list(names(disease), "disease"))
\}

meta.result <- palm(rel.abd = rel.abd, covariate.interest = covariate.interest)



\code{\LinkA{palm.get.summary}{palm.get.summary}},
\code{\LinkA{palm.null.model}{palm.null.model}},
\code{\LinkA{palm.meta.summary}{palm.meta.summary}}


Zhoujingpeng Wei <zwei74@wisc.edu>

\end{Arguments}
\HeaderA{palm.get.summary}{Generate PALM summary statistics for individual studies}{palm.get.summary}
%
\begin{Description}
This function takes the output object from the \code{palm.null.model} function and
constructs summary statistics for the covariates of interest in each study.
The resulting output can be directly used as input to the \code{palm.meta.summary} function
to perform meta-analysis.
\end{Description}
%
\begin{Usage}
\begin{verbatim}
palm.get.summary(
  null.obj,
  covariate.interest,
  cluster = NULL,
  correct = "median"
)
\end{verbatim}
\end{Usage}
%
\begin{Arguments}
\begin{ldescription}
\item[\code{null.obj}] The output object from the \code{palm.null.model} function.

\item[\code{...}] Additional arguments passed to \code{palm}.
In \code{covariate.interest} and \code{cluster}, the order of samples must match
the order in \code{rel.abd} used in the \code{palm.null.model} function.
\end{ldescription}
\end{Arguments}
%
\begin{Value}
A list containing one component per study. Each component includes the following elements:
\begin{ldescription}
\item[\code{est}] A matrix containing the association effect estimates for the study.
Rows correspond to microbial feature IDs and columns correspond to covariate IDs.
\item[\code{stderr}] A matrix containing the standard errors of the association effect estimates for the study.
Rows correspond to microbial feature IDs and columns correspond to covariate IDs.
\item[\code{n}] The sample size for the study.
\end{ldescription}
\end{Value}
%
\begin{SeeAlso}
\code{\LinkA{palm.null.model}{palm.null.model}},
\code{\LinkA{palm.meta.summary}{palm.meta.summary}},
\code{\LinkA{palm}{palm}}
\end{SeeAlso}
%
\begin{Examples}
\begin{ExampleCode}

library(PALM)
data("CRC_data", package = "PALM")
CRC_abd <- CRC_data$CRC_abd
CRC_meta <- CRC_data$CRC_meta

########## Generate summary statistics ##########
rel.abd <- list()
covariate.interest <- list()
for (d in unique(CRC_meta$Study)) {
  rel.abd[[d]] <- CRC_abd[CRC_meta$Sample_ID[CRC_meta$Study == d], ]
  disease <- as.numeric(CRC_meta$Group[CRC_meta$Study == d] == "CRC")
  names(disease) <- CRC_meta$Sample_ID[CRC_meta$Study == d]
  covariate.interest[[d]] <- matrix(disease, ncol = 1,
                                    dimnames = list(names(disease), "disease"))
}

null.obj <- palm.null.model(rel.abd = rel.abd)
summary.stats <- palm.get.summary(null.obj = null.obj,
                                  covariate.interest = covariate.interest)

\end{ExampleCode}
\end{Examples}
\HeaderA{palm.meta.summary}{Meta-analyze summary statistics across studies}{palm.meta.summary}
%
\begin{Description}
This function takes the summary statistics output from the \code{palm.get.summary} function
and combines the results across studies to perform feature-level association testing
for each covariate of interest.
\end{Description}
%
\begin{Usage}
\begin{verbatim}
palm.meta.summary(summary.stats, p.adjust.method = "fdr", meta.method = "EE")
\end{verbatim}
\end{Usage}
%
\begin{Arguments}
\begin{ldescription}
\item[\code{summary.stats}] The output object from the \code{palm.get.summary} function.

\item[\code{...}] Additional arguments passed to \code{palm}.
\end{ldescription}
\end{Arguments}
%
\begin{Value}
An object with the same structure as the output of the \code{palm} function when performing meta-analysis.
\end{Value}
%
\begin{SeeAlso}
\code{\LinkA{palm.null.model}{palm.null.model}},
\code{\LinkA{palm.get.summary}{palm.get.summary}},
\code{\LinkA{palm}{palm}}
\end{SeeAlso}
%
\begin{Examples}
\begin{ExampleCode}

library(PALM)
data("CRC_data", package = "PALM")
CRC_abd <- CRC_data$CRC_abd
CRC_meta <- CRC_data$CRC_meta

########## Generate summary statistics ##########
rel.abd <- list()
covariate.interest <- list()
for (d in unique(CRC_meta$Study)) {
  rel.abd[[d]] <- CRC_abd[CRC_meta$Sample_ID[CRC_meta$Study == d], ]
  disease <- as.numeric(CRC_meta$Group[CRC_meta$Study == d] == "CRC")
  names(disease) <- CRC_meta$Sample_ID[CRC_meta$Study == d]
  covariate.interest[[d]] <- matrix(disease, ncol = 1,
                                    dimnames = list(names(disease), "disease"))
}

null.obj <- palm.null.model(rel.abd = rel.abd)
summary.stats <- palm.get.summary(null.obj = null.obj,
                                  covariate.interest = covariate.interest)

########## Meta-analysis ##########
meta.result <- palm.meta.summary(summary.stats = summary.stats)

\end{ExampleCode}
\end{Examples}
\HeaderA{palm.null.model}{Get information from the null model}{palm.null.model}
%
\begin{Description}
The \code{palm.null.model} function computes the estimated mean microbial feature
proportions and residuals of the relative abundances under the null model for each study.
This function does not require the \code{covariate.interest} information.
The output of this function is used as input for the \code{palm.get.summary} function
to construct summary statistics for the covariates of interest in each study.
\end{Description}
%
\begin{Usage}
\begin{verbatim}
palm.null.model(
  rel.abd,
  covariate.adjust = NULL,
  depth = NULL,
  depth.filter = 0,
  prev.filter = 0.1
)
\end{verbatim}
\end{Usage}
%
\begin{Arguments}
\begin{ldescription}
\item[\code{...}] Additional arguments passed to \code{palm}.
\end{ldescription}
\end{Arguments}
%
\begin{Value}
A list containing one component per study. Each component includes the following elements:
\begin{ldescription}
\item[\code{Y\_I}] A matrix of predicted abundances.
\item[\code{Y\_R}] A matrix of abundance residuals.
\item[\code{Z}] The cleaned and formatted \code{covariate.adjust}.
\item[\code{rm.sample.idx}] Indices of the samples removed from the study.
\end{ldescription}
\end{Value}
%
\begin{SeeAlso}
\code{\LinkA{palm.get.summary}{palm.get.summary}},
\code{\LinkA{palm.meta.summary}{palm.meta.summary}},
\code{\LinkA{palm}{palm}}
\end{SeeAlso}
%
\begin{Examples}
\begin{ExampleCode}

library(PALM)
data("CRC_data", package = "PALM")
CRC_abd <- CRC_data$CRC_abd
CRC_meta <- CRC_data$CRC_meta

########## Generate summary statistics ##########
rel.abd <- list()
for (d in unique(CRC_meta$Study)) {
  rel.abd[[d]] <- CRC_abd[CRC_meta$Sample_ID[CRC_meta$Study == d], ]
}

null.obj <- palm.null.model(rel.abd = rel.abd)

\end{ExampleCode}
\end{Examples}
\printindex{}
\end{document}
